\section{Estudio de caso: Vision-guided Autonomy}
\subsection{Descomposición de tareas en escenarios reales}
La Slide 09 muestra un vision-guided manipulation case, que incluye tres módulos: grasp, transport y place; cada módulo, a su vez, anida reward, reset y goal manager autónomos. El ponente enfatiza que los escenarios reales suelen presentar long-tailed failure mode, por lo que es necesario incorporar anomaly detection en el pipeline.

\begin{figure}[H]
\centering
\includegraphics[width=0.85\textwidth]{slides-images/slide-09.jpg}
\caption{La Slide 9 muestra el flujo de trabajo de vision-guided manipulation y enfatiza cómo colaboran los distintos module.}
\end{figure}

\begin{knowledgebox}{División de Vision-guided Autonomy}
\begin{itemize}
\item Grasp module: usa visual affordance para predecir el grasp point;\\
\item Transport module: mueve el objeto mediante una goal-conditioned policy;\\
\item Place module: usa un success classifier para determinar si se completó la colocación en el objetivo.
\end{itemize}
\end{knowledgebox}

\subsection{Observabilidad y respuesta a anomalías}
El ponente sugiere registrar una anomaly metric en cada etapa: confidence drop del modelo de vision, failure rate de la reset policy y goal coverage drop. Al llevar estas métricas a un dashboard, es posible hacer trigger de un automated rollback o de una reward recalibration.

\begin{importantbox}{Matriz de respuesta a anomalías}
\begin{tabularx}{\textwidth}{lX}
\toprule
Tipo de anomalía & Respuesta \\
\midrule
Vision confidence drop & Entrar en safe mode y reducir la exploration rate; \\
Reset failure & Activar de inmediato una manual inspection y revertir al stable checkpoint más reciente; \\
Goal coverage shrink & Ampliar el goal replay buffer y volver a hacer thaw de los failed goal cluster; \\
\bottomrule
\end{tabularx}
\end{importantbox}

\subsection{Resumen de la sección}
Mediante el caso de vision-guided manipulation, se muestran las estrategias de observabilidad y de respuesta a anomalías del Autonomy pipeline en el despliegue fuera de línea y en línea.
\newpage
